\documentclass[10pt,a4paper]{amsart}
\usepackage[margin=19mm]{geometry}
\usepackage{amsmath,amssymb,mathtools}
\usepackage[scr=rsfs,cal=euler]{mathalfa}
\usepackage{xfrac}
\usepackage[unicode,hidelinks,bookmarks=false]{hyperref}
\newcommand{\ie}{i.e.\ }
\newcommand{\cf}{cf.\ }
\newcommand{\resp}{respectively\ }
\newcommand{\Subsection}{\subsection}
\providecommand{\nobreakdash}{\nobreak\hbox{-}}
\setlength{\emergencystretch}{2em}
\multlinegap0pt
\usepackage{mathtools}
\usepackage{bm}
\usepackage{amssymb}
\usepackage{float}
\usepackage{tikz}
\newtheorem{lemma}{Lemma}
\usepackage{cleveref}
\crefname{lemma}{Lemma}{Lemmas}
\title{Uncovering Symmetry Transfer in Large Language Models via Layer-Peeled Optimization}
\author{Zhehang Du, Hangfeng He, Weijie Su}
\date{Source: arXiv:2605.12756v1 (CC BY 4.0)}
\begin{document}
\maketitle

\noindent\emph{Source and licence.} Uncovering Symmetry Transfer in Large Language Models via Layer-Peeled Optimization. Version arXiv:2605.12756v1 (CC BY 4.0), distributed by arXiv under the Creative Commons Attribution 4.0 International licence, http://creativecommons.org/licenses/by/4.0/, as recorded in the arXiv licence metadata for this exact version and shown on the abstract page https://arxiv.org/abs/2605.12756v1. Full licence notice: \texttt{licenses/arxiv-2605.12756-CC-BY-4.0.txt}.\par\smallskip

\noindent\textit{Editorial scope.} Counted unit: the article's lemma lem:softmax\_ce\_strict\_convexity\_V0 (source0.tex lines 1807--1821). The article's complete original proof of it is delivered with it (source0.tex lines 1823--1865, one \texttt{proof} environment of 43 lines). No mathematics is written, repaired or re-derived: the statement and the proof are the article's own lines, in the article's own notation, and no retained line is substituted or re-flowed.  No further source-local context is needed: every cross-reference of every retained line resolves inside the two delivered spans.  Nothing was omitted from any retained span, so the package carries no omission record.  No retained line contains a citation, so the wrapper carries no bibliography and no bibliography entry.  No retained line contains a figure environment, an image-inclusion command or drawing code, so no image is delivered and the package declares no asset dependency.  Editorial formatting only, in addition: an AMS \texttt{amsart} preamble that restates the article's own declarations and macros used by the retained lines, namely the environments \texttt{lemma} on separate counters, together with the standard packages those lines need.  The retained environments print in this document as Lemma 1 in order of appearance, whereas the article prints the same items with the numbering of its own full document.  The complete native source of the article as distributed in the arXiv e-print for this exact version is bundled unchanged as \texttt{sources/200446/source0.tex}.
\begin{lemma}[Strict Convexity of Softmax Cross-Entropy on the Zero-Sum Subspace]
\label{lem:softmax_ce_strict_convexity_V0}
Let \(f(\boldsymbol u) \coloneqq \mathcal L(\sigma(\boldsymbol u), \boldsymbol p)\)
with \(\boldsymbol p\in\Delta^{m-1}\) fixed and \(\boldsymbol q \coloneqq \sigma(\boldsymbol u)\).
Then \(f\) is convex on \(\mathbb R^m\) with Hessian
\[
\nabla^2 f(\boldsymbol u)   =   \mathrm{diag}(\boldsymbol q) - \boldsymbol q \boldsymbol q^\top \succeq 0.
\]
Moreover, \(f\) is strictly convex along every direction in the zero-sum subspace
\[
\mathcal V_0 \coloneqq \{\boldsymbol v\in\mathbb R^m : \boldsymbol{1}_m^\top \boldsymbol v = 0 \}.
\]
Equivalently, for every nonzero \(\boldsymbol v\in \mathcal V_0\),
\(\boldsymbol v^\top \nabla^2 f(\boldsymbol u)   \boldsymbol v > 0\).
\end{lemma}

\begin{proof}[Proof of~\cref{lem:softmax_ce_strict_convexity_V0}]
First, we compute the gradient of \(f(\boldsymbol{u})\),
\[ \frac{\partial f}{\partial u_j} = \frac{\partial}{\partial u_j} \left( \log\left(\sum_{k=1}^m e^{u_k}\right) \right) - \frac{\partial}{\partial u_j} \left( \sum_{k=1}^m p_k u_k \right). \]
The derivative of the log-sum-exp term is
\[ \frac{\partial}{\partial u_j} \log\left(\sum_{k=1}^m e^{u_k}\right) = \frac{1}{\sum_{k=1}^m e^{u_k}} \cdot e^{u_j} = (\sigma(\boldsymbol{u}))_j. \]
The derivative of the linear term is simply \(p_j\).
Letting \(\boldsymbol{q} = \sigma(\boldsymbol{u})\), the gradient vector is
\[ \nabla_{\boldsymbol{u}} f(\boldsymbol{u}) = \boldsymbol{q} - \boldsymbol{p}. \]
Next, we compute the Hessian matrix, \(H_f = \nabla^2_{\boldsymbol{u}} f\), whose entries are \( (H_f)_{ij} = \partial^2 f/(\partial u_i \partial u_j) \). This is equivalent to finding the Jacobian matrix of the gradient \(\nabla_{\boldsymbol{u}} f = \boldsymbol{q} - \boldsymbol{p}\). Since \(\boldsymbol{p}\) is a constant vector, the Hessian is simply the Jacobian of \(\boldsymbol{q} = \sigma(\boldsymbol{u})\).
The partial derivative of a component \(q_j\) with respect to a logit \(u_i\) is
\[ \frac{\partial q_j}{\partial u_i} = \frac{\partial}{\partial u_i}\left( \frac{e^{u_j}}{\sum_k e^{u_k}} \right). \]
If \(i = j\),
    \[ \frac{\partial q_j}{\partial u_j} = \frac{e^{u_j}\left(\sum_k e^{u_k}\right) - e^{u_j}(e^{u_j})}{\left(\sum_k e^{u_k}\right)^2} = \frac{e^{u_j}}{\sum_k e^{u_k}} - \left(\frac{e^{u_j}}{\sum_k e^{u_k}}\right)^2 = q_j - q_j^2 = q_j(1-q_j). \]
If \(i \neq j\),
    \[ \frac{\partial q_j}{\partial u_i} = \frac{0 \cdot (\sum_k e^{u_k}) - e^{u_j}(e^{u_i})}{\left(\sum_k e^{u_k}\right)^2} = -\frac{e^{u_j}}{\sum_k e^{u_k}}\frac{e^{u_i}}{\sum_k e^{u_k}} = -q_i q_j. \]
Combining these results, the Hessian matrix is
\[ H_f = \nabla^2_{\boldsymbol{u}} f = \mathrm{diag}(\boldsymbol{q}) - \boldsymbol{q}\boldsymbol{q}^\top. \]
To prove that \(f(\boldsymbol{u})\) is convex, it suffices to show that \(H_f\) is positive semidefinite. For any \(\boldsymbol{v} \in \mathbb{R}^m\),
\[
\boldsymbol{v}^\top H_f \boldsymbol{v} = \boldsymbol{v}^\top (\mathrm{diag}(\boldsymbol{q}) - \boldsymbol{q}\boldsymbol{q}^\top) \boldsymbol{v} = \boldsymbol{v}^\top \mathrm{diag}(\boldsymbol{q}) \boldsymbol{v} - \boldsymbol{v}^\top \boldsymbol{q}\boldsymbol{q}^\top \boldsymbol{v}.
\]
Expanding both terms gives
\[
\boldsymbol{v}^\top H_f \boldsymbol{v} = \sum_{i=1}^m q_i v_i^2 - \left(\sum_{i=1}^m q_i v_i\right)^2.
\]
Thus \(\boldsymbol{v}^\top H_f \boldsymbol{v}\) is exactly the variance of the discrete random variable taking values \(\{v_1,\dots,v_m\}\) with weights \(\{q_1,\dots,q_m\}\):
\[
\boldsymbol{v}^\top H_f \boldsymbol{v} = \mathbb{E}[v^2] - (\mathbb{E}[v])^2 = \mathrm{Var}_q(v).
\]
Equivalently,
\[
\boldsymbol{v}^\top H_f \boldsymbol{v} = \mathrm{Var}_q(v) = \mathbb{E} \left[(v-\mathbb{E}[v])^2\right] \ge 0.
\]
Hence \(H_f \succeq 0\), so \(f\) is convex on \(\mathbb{R}^m\).

Now let \(\boldsymbol{v}\in \mathcal V_0\) be nonzero. To prove strict convexity on \(\mathcal V_0\), it remains to show that \(\mathrm{Var}_q(v) > 0\). Since \(\boldsymbol{q} = \sigma(\boldsymbol{u})\), every coordinate satisfies \(q_i > 0\). Because \(q_i>0\) for every \(i\), \(\mathrm{Var}_q(v)=0\) holds if and only if all coordinates \(v_i\) are equal. Therefore
\[ \boldsymbol{v}^\top H_f \boldsymbol{v} = 0 \quad \iff \quad v_1 = v_2 = \cdots = v_m. \]
Thus \(\boldsymbol{v}^\top H_f \boldsymbol{v} = 0\) if and only if \(\boldsymbol{v}\) is a multiple of \(\boldsymbol{1}_m\). But a nonzero vector in
\[
\mathcal V_0 = \{\boldsymbol{v} \in \mathbb{R}^m : \boldsymbol{1}_m^\top \boldsymbol{v} = 0\}
\]
cannot be a multiple of \(\boldsymbol{1}_m\). Therefore \(\boldsymbol{v}^\top H_f \boldsymbol{v} > 0\) for every nonzero \(\boldsymbol{v} \in \mathcal V_0\). Hence the Hessian is positive definite on \(\mathcal V_0\), and \(f\) is strictly convex along that subspace.
\end{proof}

\end{document}
